\begin{statement}[C.4.3; Proved: algebraicity and explicit fields for every order]
Let $t_0\in\mathbb C\setminus\{0,1\}$ and suppose $P_{t_0}$ has
exact order $n\ge2$. Then $t_0$ and $N(P_{t_0})$ are algebraic.
Put $K_0=\mathbb Q(t_0,\bar t_0)$,
\[
 W=\Delta(t_0)\Delta(\bar t_0)>0,\qquad
 V=\psi_{n+1}(P_{t_0})\psi_{n+1}(P_{\bar t_0})>0.
\]
With positive real radicals,
\[
 N(P_{t_0})^{12n}=W^{n/2}V^{-6/(n+2)}
 \ \in\ K_0\bigl(W^{1/2},V^{1/(n+2)}\bigr).
\]
This supplies an explicit field, without a claim that it is minimal.
More generally, if $m\equiv\pm1\pmod n$, then
\[
 \lambda_\infty(P_{t_0})
 =\frac{\log|\psi_m(P_{t_0})|}{m^2-1}
       -\frac1{12}\log|\Delta(t_0)|.
\]
These formulas determine the height and $N^{12}$ on the entire
torsion locus.
\end{statement}
\begin{proof}
The division polynomial $\psi_n(P_t)$ is a nonzero rational
polynomial in $t$. To check nonvanishing for every $n$, a hypothetical
torsion identity would express the elliptic logarithm as
$a\tau+b$ with constant rational $a,b$. Its bounded cusp limit
forces $a=0$, but that limit has nonzero imaginary part, a
contradiction. Thus $t_0$ is a zero of a nonzero polynomial over
$\mathbb Q$. Formula \eqref{eq:cs-torsionvalue} shows algebraicity
of $N$ using only positive rational powers of nonzero algebraic
absolute values. Taking $m=n+1$, raising to the $n$th power, and
using $|\Delta|=W^{1/2}$, $|\psi_{n+1}|=V^{1/2}$ gives the exact
field statement. The local-height formula is its logarithm.

Here is a polynomial relation involving only $t$ and $r=N^{12}$,
if one prefers that form to positive radicals. Remove the factors
at $0,1$ from $\psi_n(P_t)$ and divide successively by the
corresponding polynomials for proper divisors of $n$; let $F_n$ be
the resulting primitive polynomial. At a point of exact order
$d\mid n$, the formal group gives
$z([n/d]Q)=(n/d)z(Q)+O(z(Q)^2)$, so the vanishing multiplicities for
$\psi_d$ and $\psi_n$ agree. Hence this cancellation leaves exactly
the exact-order-$n$ divisor, with its multiplicities. Set
$m=n+1$ and $e=m^2-1$. Clearing rational denominators in
\[
 \operatorname{Res}_S\left(
 F_n(S),\
 r^{2e}\{\psi_m(P_T)\psi_m(P_S)\}^{12}
              -\{\Delta(T)\Delta(S)\}^{e}\right)
\]
gives a polynomial in $\mathbb Q[T,r]$ that vanishes at
$(t_0,N(P_{t_0})^{12})$, since $S=\bar t_0$ is a common root.
It is nonzero: at an exact-order-$n$ root no $\psi_m$ vanishes,
so the leading coefficient in $r$ of the resultant is nonzero.
This is an explicit recurrence-and-resultant prescription for
every $n$, including nonreal torsion parameters.
\end{proof}
